\newcommand{\R}{\mathbb{R}}
\newcommand{\Z}{\mathbb{Z}}
\newcommand{\Q}{\mathbb{Q}}
\newcommand{\K}{\mathbb{K}}

\newcommand{\vectorComponents}[1]{#1_1, \dots, #1_n}
\newcommand{\ideal}[1]{\langle #1 \rangle}
\newcommand{\vectorComponentsX}[2]{#1_1, \dots, #1_#2}
\newcommand{\xalpha}{x^\alpha}
\newcommand{\polyRing}{k\lbrack \vectorComponentsX{x}{n} \rbrack}
\newcommand{\finiteIdeal}[3]{#1(#2_1), \dots, #1(#2_#3)}
\newcommand{\specialGrobCondition}{\ideal{LT(I)} = \ideal{\finiteIdeal{LT}{g}{t}}}
\newcommand{\sPoly}[2]{\overline{#1}^{#2}}
\setpapersize{A4}
\setmarginsrb{30mm}{20mm}{30mm}{30mm}%left top right bottom
                 {8mm}{6mm}{10mm}{10mm}%headheight headsep footheight footskip

% Some useful colors
\definecolor{Rosso}{RGB}{232,61,61}
\definecolor{Viola}{RGB}{188,30,188}
\definecolor{Celeste}{RGB}{34,186,211}
\definecolor{Arancione}{RGB}{232,181,19}
\definecolor{Blu}{RGB}{0, 44, 92}
\definecolor{Azzurro}{RGB}{57,186,238}


\hypersetup{colorlinks=true,linkcolor=Blu,allbordercolors=white, urlcolor=blue, citecolor=Viola}
\usepackage{enumitem} 
\renewcommand\labelitemi{\textbullet} 


% math environments
%  we use \newtheorem to specify our theorem-like environments
   % (theorem, definition, etc.) and how to display and number them.
   %
   % Note: The \theoremstyle declarations affect the appearance of the
   % Theorems, Definitions, etc.
   % Note 2: If you want swedish names you have to go and change the argument in brackets e.g. \newtheorem{thm}[defi]{Theorem} should become \newtheorem{thm}[defi]{Sats}
\theoremstyle{definition}
\newtheorem{definition}{Definition}[section]

\newtheorem*{problem}{Problem}
\newtheorem{conjecture}[definition]{Conjecture}

\theoremstyle{plain}
\newtheorem{theorem}[definition]{Theorem}
\newtheorem{proposition}[definition]{Proposition}
\newtheorem{corollary}[definition]{Corollary}
\newtheorem{lemma}[definition]{Lemma}

\theoremstyle{remark}
\newtheorem*{question}{Question}
\newtheorem{remark}[definition]{Remark}
\newtheorem{example}[definition]{Example}
\newtheorem*{notation}{Notation}

\onehalfspacing
\usepackage{algorithm}
\usepackage{algpseudocode}